% !TEX root = ../main.tex
% Chapter 1
\chapter{Introducción} % Main chapter title
\label{sec:Introduccion} % For referencing the chapter elsewhere, use


La última milla constituye el eslabón final y más crítico de la cadena de suministro, encargada de conectar de manera directa el producto de una empresa con el consumidor privado final \cite{boysen2020last}. Asimismo, en el escenario contemporáneo, fuertemente impulsado por el auge del comercio electrónico y los servicios de entrega ultra-rápida, este proceso ha adquirido una relevancia estratégica debido a su alta complejidad operacional y sus elevados costos asociados. Así, diversos estudios en el ámbito de la ingeniería de transporte y la logística urbana señalan que la última milla es típicamente considerada la parte más ineficiente de la cadena, llegando a representar cerca del 53\% de los costos logísticos totales que enfrentan las organizaciones \cite{yadav2025assessing}. 

Por consiguiente, para las empresas del sector retail y de distribución es un imperativo crucial manejar eficientemente sus recursos. De manera que, la optimización nace como una herramienta fundamental para mejorar el manejo de recursos, reducir el consumo y organizar flotas de vehículos en tiempos viables, una tarea que sobrepasa las capacidades de planificación puramente manuales de un ser humano.

\subsection{Descripción del Entorno Operacional: Simpliroute}
Ahora bien, la oportunidad de esta investigación se genera y contextualiza en el entorno operativo de Simpliroute, una plataforma comercial de tecnología para la logística dedicada a ayudar a empresas ofreciendo entregas rápidas, eficientes e inteligentes.

Así, los clientes de Simpliroute, que incluyen empresas de retail, farmacias y restaurantes, utilizan la plataforma para optimizar sus operaciones de entrega en la última milla. De manera que, la plataforma permite a los usuarios ingresar datos de pedidos, ubicaciones de entrega y restricciones de tiempo, generando rutas optimizadas para los conductores. 

Sin embargo, se ha realizado un análisis para verificar el seguimiento de las rutas generadas por Simpliroute. Donde, se ha notado que sus conductores ignoran las recomendaciones y deciden tomar su propia ruta por razones que no se han identificado, notandose una brecha entre la ruta óptima generada por el algoritmo y la ruta ejecutada en el mundo real. 

En consecuencia, este fenómeno de no adherencia a la ruta óptima representa un desafío crítico para la eficiencia operativa de Simpliroute, ya que impacta directamente en el uso de recursos de los clientes e incluso en los costos logísticos y en la satisfacción del cliente.

